\scanpage{215}
\margin{p.d. dans les anneaux \unsure{alternés}}
% margin note written sideways
\margin{Feuille paumée\ldots}
\noindent$\left[\vphantom{\rule{0pt}{6\baselineskip}}\right.$\begin{minipage}{.90\linewidth}
\uline{Théorème} Soient $S$ un anneau muni d'un idéal $I$ à puiss. divisées, $M^\ast=M^0+M^1$ un $S$-module gradué mod $2$,
$A(M^{\deleted{\unsure{0}}\ast})=\Gamma^\ast(M^0)\otimes_S\bigwedge^\ast M^1$ \inserted{$S$-algèbre} gradué par la parité du degré total, [et augmentée vers $S$ de la façon évidente $\varepsilon:A\longrightarrow S$]
$J=$ l'idéal engendré dans $A$ par $I$ ($\subset S\subset A$) et par $M^\ast$ ($\subset S$), [donc $J=\varepsilon^{-1}(I)$].
\end{minipage}

a) On peut trouver sur \deleted{$A$} $J\subset A$ une structure à p.d. telle que 1°) $S\longrightarrow A$ soit un hom. d'anneaux compatible aux p.d. sur $I\subset S$ et $J\subset A$ \quad 2°) Pour $(S,M^\ast)$ variable, cette structure à p.d. est fonctorielle. Ces conditions caractérisent de façon unique la structure à p.d. envisagée.

b) Soit $(B^\ast\deleted{\unsure{,K^\ast}})$ une $S$-algèbre graduée mod $2$, \deleted{alternée strict} alternée, munie d'un idéal $K^\ast$ gradué à p.d., \deleted{l'application de restriction} \inserted{compatibles avec les p.d.} sur $I\subset S$.
L'application de restriction
\[
\begin{aligned}
&\operatorname{Hom}_{S\text{-alg à p.d.}}\bigl((A(M^\ast),J(M^\ast)),(B^\ast,K^\ast)\bigr)\\
&\hspace{8em}\longrightarrow\operatorname{Hom}_{S\text{-mod grd}}(M^\ast,K^\ast)
\end{aligned}
\]
est bijective. (Donc il y a lieu d'appeler $A(M^\ast)$
la \unsure{$B$|$S$}-algèbre \inserted{graduée alternée} à p.d. \deleted{sur} $S$ \uline{engendrée} par le $S$-module gradué $M^\ast$.)

\uline{Dém} :
% archive sheet number: 214
